% ============================================================================
% 02_literature.tex — Related literature
% ============================================================================

\section{Related literature}
\label{sec:literature}

The paper draws on, and speaks to, five strands of work.

\emph{User cost and tenure choice.} The theoretical benchmark for rent-versus-own comparisons is the user cost of owner-occupied housing --- foregone interest on equity, mortgage interest, maintenance and depreciation, less expected capital gains --- formalised by \textcite{poterba1984} on tenure-choice foundations laid by \textcite{henderson1983}, and turned into the standard housing-valuation tool by \textcite{himmelberg2005}. A measurement branch shows, however, that ex-ante user costs diverge persistently from observed rents and are acutely sensitive to imputed appreciation and financing assumptions \parencite{verbrugge2008,diaz2008}, and that credible price-rent comparisons require quality- and composition-matched data \parencite{hill2016,bracke2015}. This paper is the ex-post counterpart to that ex-ante tradition: rather than imputing an equilibrium user cost, I track the realised cash flows of both tenures month by month, which sidesteps the imputation problem; my required appreciation rate is the user-cost equilibrium condition solved for the capital-gains term rather than assuming it; and the composition-consistent national aggregate of Section~\ref{sec:representative} responds directly to the matched-measurement critique --- the official England price and rent aggregates, I show, fail exactly that test.

\emph{Buy-versus-rent horse races.} The closest methodological ancestors are the ex-post backtests for the United States: \textcite{beracha2012}, who run cash-flow-matched buy-versus-rent races over 1978--2009 and find renting frequently competitive; the companion indifference-appreciation exercise of \textcite{berachaseiler2012}, an antecedent of my RAR metric computed under stylised ex-ante assumptions rather than realised paths; and \textcite{rappaport2010} and \textcite{goodman2018} on homeownership as wealth-building, with \textcite{herbert2013} surveying the distributional evidence. For Britain, the only comparable academic exercise is \textcite{mostafa2019}, an annual-frequency comparison of first-time-buyer cohorts over 1975--2011 in which the renter's counterfactual earns building-society deposit rates; buying wins essentially always. \textcite{tabner2016} provides a UK net-present-value framework under assumed parameters, and practitioner analyses exist for Canada \parencite{felix2025}. Relative to this work, the present paper is --- to my knowledge --- the first monthly, cash-flow-matched backtest on UK data, and the first in which the British renter's counterfactual is a diversified equity portfolio; that single change overturns the always-buy conclusion and produces near-parity, while extending the sample through the post-2012 boom, the pandemic, and the 2022--2023 rate shock that earlier UK samples miss.

\emph{Housing returns.} A recent measurement literature quantifies long-run returns to residential property: house prices track incomes rather than compounding like equities over the long run \parencite{knoll2017}; total returns are competitive with equities only once rental yield is included and before costs \parencite{jorda2019rore}; and micro-level studies show that transaction costs, management and idiosyncratic risk substantially erode net returns to individual owners \parencite{chambers2021,eichholtz2021total,giacoletti2021,demers2022}, with large spatial heterogeneity \parencite{amaral2025}. My results are the household-level, England-specific counterpart: with realistic costs the unlevered housing return loses badly to equities --- an all-cash buyer never wins --- and it is mortgage leverage, not the housing return itself, that makes ownership financially competitive.

\emph{Housing in household portfolios.} A structural literature examines how owner-occupied housing shapes household portfolios and lifecycle choices \parencite{flavin2002,cocco2005,yao2005,chetty2017,vestman2019}; \textcite{campbell2006} and \textcite{gomes2021} survey the field. These models emphasise precisely the two forces my backtest measures ex post: the portfolio distortion of a large, levered, indivisible housing position, and the financing channel through which mortgage debt amplifies returns to home equity. The empirical leverage gradient I document --- buying wins two-thirds of cohorts at 95\% LTV and none at all in cash --- is, in effect, a realised-data trace of the mechanism these models study.

\emph{Rent risk, commitment, and the UK market.} Two channels are deliberately switched off in my accounting and frame its interpretation. Ownership hedges rent risk \parencite{sinai2005}, while my renter bears actual market rents throughout; and mortgage amortisation acts as forced saving for present-biased households \parencite{laibson1997,thaler2004,schlafmann2021}, with \textcite{bernstein2024} showing amortisation causally builds net worth and \textcite{sodini2023} identifying substantial benefits of ownership quasi-experimentally --- whereas my renter is endowed with perfect investment discipline. Near-parity for a perfectly disciplined renter therefore implies an advantage for ownership among typical households, which is how I read the result. Finally, the paper connects to UK housing economics --- cycles and regional dynamics \parencite{muellbauer1997,meen1999}, rate-driven valuations \parencite{milesmonro2021}, supply constraints \parencite{hilbervermeulen2016}, and the stamp-duty responses documented by \textcite{besley2014} and \textcite{bestkleven2018}, which motivate my period-accurate SDLT module --- and to the evidence that price-rent ratios forecast housing returns \parencite{caseshiller1989,campbellshiller1988,gallin2008,campbell2009,engstedpedersen2015}: my finding that the entry price-to-rent ratio and mortgage rate explain 74\% of cohort outcomes is the cash-flow-matched analogue of that predictability.
